% !TeX program = lualatex
% !TeX encoding = UTF-8
\documentclass[11pt,a4paper]{ltjsarticle}
\usepackage{amsmath,amssymb,amsthm,amsfonts,mathtools}
\usepackage{lmodern}
\usepackage{geometry}
\usepackage{enumitem}
\usepackage{xcolor}
\usepackage{booktabs}
\usepackage{tabularx}
\usepackage{array}
\usepackage{hyperref}
\usepackage{tikz}
\usetikzlibrary{cd,positioning,arrows.meta}
\geometry{margin=1in}
\setlength{\emergencystretch}{3em}
\raggedbottom

%-------------------------------------------------
%  定理環境
%-------------------------------------------------
\theoremstyle{definition}
\newtheorem{definition}{定義}[section]
\newtheorem{axiom}{公理}[section]
\newtheorem{lemma}{補題}[section]
\newtheorem{theorem}{定理}[section]
\newtheorem{corollary}{系}[section]
\newtheorem{proposition}{命題}[section]
\newtheorem{remark}{注記}[section]
\newtheorem{example}{例}[section]

\renewcommand{\proofname}{\textbf{証明.}}
\renewcommand{\abstractname}{Abstract}

%-------------------------------------------------
%  マクロ
%-------------------------------------------------
\newcommand{\Sk}{\mathbf{Sk}}
\newcommand{\Cat}{\mathbf{Cat}}
\newcommand{\Mnd}{\mathbf{Mnd}}
\newcommand{\End}{\mathbf{End}}
\newcommand{\id}{\mathrm{id}}
\newcommand{\Hom}{\mathrm{Hom}}
\newcommand{\Nat}{\mathrm{Nat}}
\newcommand{\Asm}{\mathrm{Asm}}
\newcommand{\Real}{\Vdash}

% 五蘊モナド
\newcommand{\Mdelta}{\mathcal{M}_{\delta}}
\newcommand{\Mneg}{\mathcal{M}_{\neg\neg}}
\newcommand{\Mr}{\mathcal{M}_{\mathrm{r}}}
\newcommand{\Mc}{\mathcal{C}}
\newcommand{\T}{\mathcal{T}}

% 随伴・スパイラル
\newcommand{\Spiral}{\mathcal{S}}
\newcommand{\SpiralT}{\mathcal{S}_{\mathcal{T}}}
\newcommand{\fold}{\mu}
\newcommand{\unfold}{\delta}
\newcommand{\trace}{\mathrm{Tr}}

% その他
\newcommand{\ext}{\mathrm{ext}}
\newcommand{\intn}{\mathrm{int}}
\newcommand{\TExt}{T_{\ext}}
\newcommand{\TInt}{T_{\intn}}

% E_∞関連
\newcommand{\PrT}{\mathrm{Pr}_{\mathcal{T}}}
\newcommand{\PrSk}{\mathrm{Pr}_{\Sk}}
\newcommand{\LawvereInfinity}{\Lambda_\infty}
\newcommand{\SigmaT}{\Sigma_{\mathcal{T}}}
\newcommand{\SigmaSk}{\Sigma_{\Sk}}
\newcommand{\Skobs}{\mathop{\mathrm{Sk}}_{\mathrm{obs}}}
\newcommand{\Skiso}{\mathop{\mathrm{Sk}}_{\mathrm{iso}}}
\newcommand{\trunc}[1]{\left\lVert #1\right\rVert_{-1}}
\newcommand{\mweq}{\mathrel{\sim}}
\newcommand{\Kstab}{\mathcal{K}^{\mathrm{st}}}
\newcommand{\Kobs}{\mathcal{K}_{\mathrm{obs}}}
\newcommand{\KSfour}{\mathcal{K}_{S_4}}
\newcommand{\Residual}{\mathsf{C}}

\DeclareMathOperator*{\bigsqcap}{\sqcap}

%-------------------------------------------------
\begin{document}

\title{五蘊圏（Skandha Category）：\\
観測局所化による外点$\bigcirc$と統制構造}
\author{千葉将臣}
\date{2026-07-17}
\maketitle

\begin{abstract}
本稿は五蘊圏（Skandha Category）$\Sk$を圏論的に公理化し、従来
「空」または外点と呼ばれていた$\bigcirc$へ観測局所化による数学的定義を与える。
五蘊圏は、5つのモナド（搾取・二重否定・洗練・カット除去・統制）を
随伴・自然変換・分配律のネットワークとして統合する圏論的構造である。
統制モナド $\T$ は固定点方程式 $\T \cong F(\T)$ により定義される。
外点の定義には五蘊圏の安定$\infty$圏的実現$\Kstab$と完全局所化
$L:\Kstab\to\Kobs$を用い、不可視残余を
$\bigcirc_X:=\operatorname{fib}(X\to iLX)$、$S_4$領域を$\ker L$として定める。
これにより$S_3$の観測可能成分と$S_4$の不可視残余は、同一対象から分解される
相補的成分として記述される。Lawvere無限階層$\LawvereInfinity$は証明論的
自己言及固定点として別個の対象とし、$\bigcirc$との同一視を行わない。
本稿は五蘊圏を2圏論・ホップf代数・トレース付きモノイダル圏の
交差点として位置づけ、仏教哲学の五蘊・四句分別・二諦と
圏論的構造の間の同型対応を明示する。
\end{abstract}

\tableofcontents
\newpage

%=============================================================
\section{序論：五蘊圏の動機と外点の問題}
%=============================================================

\subsection{問題の設定}

圏論において、モナドは計算の文脈を記述する基本的な構造である。
Moggi \cite{Moggi1991} 以来、副作用・状態・継続・例外などは
モナドの枠組みで統一的に扱われてきた。

一方、複数のモナドが相互作用する体系において、
個別のモナドの性質だけでなく、それらの間の\textbf{分配律}・\textbf{合成}・\textbf{残余}を
記述する枠組みが必要となる。

本稿は、以下の問いに答えることを目的とする。

\begin{enumerate}
\item 5つのモナドを統合する圏論的構造はどのように公理化されるか。
\item その構造は、既存の高次圏論（2圏論、ホップf代数、トレース構造）と
  どのように接続されるか。
\item 仏教哲学の五蘊・四句分別・二諦は、圏論のどの構造に対応するか。
\item 従来「空」と呼ばれてきた$\bigcirc$を、観測局所化の不可視残余として
  どのように内在的に定義できるか。
\end{enumerate}

\subsection{五蘊圏の直観}

五蘊圏 $\Sk$ は、以下の直観に基づく。

\begin{itemize}
\item \textbf{色}（形・物質）：対象の外部構造（搾取モナド $\Mdelta$）
\item \textbf{受}（感受）：対象の安定化（二重否定モナド $\Mneg$）
\item \textbf{想}（知覚・表象）：対象の前進変換（洗練モナド $\Mr$）
\item \textbf{行}（意志・形成力）：対象の正規化（カット除去モナド $\Mc$）
\item \textbf{識}（識別・意識）：対象の自己言及的統制（統制モナド $\T$）
\end{itemize}

識は、他の4蘊を統合しつつ自身も空であるという
仏教的構造を、圏論的固定点 $\T \cong F(\T)$ として実現する。

\subsection{外点と\texorpdfstring{$\LawvereInfinity$}{Lambda-infinity}の分離}

本稿では外点$\bigcirc$とLawvere無限階層$\LawvereInfinity$
\cite{Chiba2026limit}を別個の数学的対象として扱う。前者は観測局所化の核に属する
圏論的対象、後者は対角補題による証明論的固定点であり、両者は異なる普遍問題の解である。

\begin{table}[htbp]
\centering
\caption{外点$\bigcirc$と$\LawvereInfinity$の分離}
\small
\begin{tabularx}{\textwidth}{@{}l >{\raggedright\arraybackslash}X >{\raggedright\arraybackslash}X@{}}
\toprule
\textbf{側面} & \textbf{外点$\bigcirc$} & \textbf{$\LawvereInfinity$} \\
\midrule
定義の文脈 & 観測局所化・ファイバー & 証明論・対角補題 \\
数学的型 & $\ker L$の対象または生成元 & 自己言及不動点を表す文 \\
存在条件 & 局所化の核の非自明性 & 対角補題の仮定 \\
役割 & 観測によって消去される残余 & 自己言及によって固定される命題 \\
現段階の関係 & \multicolumn{2}{l}{同一視しない。翻訳関手による接続は将来課題とする。} \\
\bottomrule
\end{tabularx}
\end{table}

この分離により、五蘊圏は$\LawvereInfinity$に依存せず、観測局所化を備えた
圏論的構造として自立する。

%=============================================================
\section{五蘊圏の基本構造}
%=============================================================

\subsection{五蘊圏の基本定義}

\begin{definition}[五蘊圏 $\Sk$]\label{def:skandha}
五蘊圏（Skandha Category）は、以下のデータからなる圏（あるいは2圏）である。
\begin{enumerate}[label=(\arabic*)]
\item \textbf{基礎圏}：対象を持つ圏 $\mathcal{C}$（通常はトポスまたは実現子論的圏 $\Asm(\mathcal{K})$）。
\item \textbf{五蘊モナド}：$\mathcal{C}$ 上の5つの自己関手
  \begin{equation}
  \Mdelta,\; \Mneg,\; \Mr,\; \Mc,\; \T : \mathcal{C} \to \mathcal{C}
  \end{equation}
  およびそれぞれに付随するモナド構造
  $(\Mdelta, \eta^\delta, \mu^\delta)$,
  $(\Mneg, \eta^\neg, \mu^\neg)$,
  $(\Mr, \eta^r, \mu^r)$,
  $(\Mc, \eta^c, \mu^c)$,
  $(\T, \eta^\T, \mu^\T)$。
\item \textbf{統制固定点}：統制モナド $\T$ は以下の固定点方程式を満たす。
  \begin{equation}
  \T \cong F(\T), \qquad F(X) = X(\Mdelta, \Mneg, \Mr, \Mc, X)
  \label{eq:fixed-point}
  \end{equation}
  ここで $F$ は4つのモナドと $X$ 自身を引数に取る関手である。
\item \textbf{識軸スパイラル}：統制モナド $\T$ による軸的作用
  \begin{equation}
  \SpiralT : (\Mdelta, \Mneg, \Mr, \Mc) \longrightarrow 
             (\Mneg, \Mr, \Mc, \Mdelta)
  \label{eq:spiral}
  \end{equation}
  これは4つのモナドを巡回する自然変換の列である。
  4回の反復は観測局所化の後で同一化される：
  \begin{equation}
  L\SpiralT^4 \simeq L.
  \end{equation}
  局所化前のずれは不可視残余関手$\Residual$によって測る。
\item \textbf{フロベニウス条件}：前進（折り畳み）$\fold^\T : F(\T) \to \T$ と
  後退（展開）$\unfold^\T : \T \to F(\T)$ の間に
  \begin{equation}
  (\fold^\T \circ \T\unfold^\T) \cong (\unfold^\T \circ \fold^\T)
  \label{eq:frobenius}
  \end{equation}
  が成立する。
\item \textbf{二諦の公理}：外在 $\TExt$ と内在 $\TInt$ の二重構造が存在し、
  相互に完全に観測不可能である。
  \begin{equation}
  \TExt \not\leftrightarrow \TInt, \qquad
  \TInt \not\leftrightarrow \TExt.
  \end{equation}
\end{enumerate}
\end{definition}

\begin{remark}
五蘊圏は通常の圏を超え、2圏的または準トポス的な構造を持つ。
特に、統制モナドの固定点方程式 \eqref{eq:fixed-point} は
圏 $\mathcal{C}$ が十分な極限を持つことを要請する。
\end{remark}

%=============================================================
\section{五句計算論基盤と観測局所化}
\label{sec:penta-foundation}
%=============================================================

\subsection{五蘊圏の五句計算論的型付け}

\begin{definition}[五句型付き五蘊圏]
五蘊圏$\Sk$に、表\ref{tab:penta-typing}の五句計算論的型付けを与える。
この型付けでは、圏論的構造を五句段階の代替物と同一視するのではなく、
各構造がどの記述可能性段階で観測されるかを指定する。
\end{definition}

\begin{table}[htbp]
\centering
\caption{五蘊圏の五句計算論的型付け}
\label{tab:penta-typing}
\small
\begin{tabularx}{\textwidth}{@{}>{\raggedright\arraybackslash}X >{\raggedright\arraybackslash}X c@{}}
\toprule
圏論的構造 & 五句計算論的型 & 段階 \\
\midrule
対象$X\in\mathcal{K}$ & $\operatorname{hlevel}(X)\geq0$ & $S_1$ \\
射$f:X\to Y$ & $\operatorname{hlevel}(f)\geq-1$ & $S_2$ \\
同型類$[X]_{\cong}$ & $\trunc{X}$ & $S_3$ \\
自然変換$\alpha:F\Rightarrow G$ & $\operatorname{Nat}(F,G)$ & $S_1$ \\
観測局所化$L$ & $L:\Kstab\to\Kobs$ & $S_3$ \\
不可視残余$\Residual(X)$ & $\operatorname{fib}(X\to iLX)$ & $S_4$ \\
\bottomrule
\end{tabularx}
\end{table}

\begin{definition}[五句truncationと観測局所化]
命題切捨てを
\[
 \trunc{\mathord{\cdot}}:\mathcal{K}\longrightarrow\trunc{\mathcal{K}},
 \qquad \trunc{\mathord{\cdot}}:S_1\longrightarrow S_3
\]
と型付けする。外点の定義には、この観測操作の安定$\infty$圏的実現として
以下の完全局所化を用いる。
\end{definition}

\subsection{安定実現と不可視残余}

\begin{definition}[観測局所化]
五蘊圏の\textbf{安定実現}を、零対象と有限極限・有限余極限を持つ安定$\infty$圏
$\Kstab$とする\cite{LurieHA}。観測局所化とは、完全関手と充満忠実な右随伴
\[
 L:\Kstab\rightleftarrows\Kobs:i,
 \qquad L\dashv i,
\]
および単位$\eta:\id_{\Kstab}\Rightarrow iL$からなる。$\Kobs$を観測可能骨格と呼ぶ。
\end{definition}

\begin{definition}[不可視残余]
各$X\in\Kstab$に対し、不可視残余関手$\Residual$と外点候補$\bigcirc_X$を
\[
 \Residual:=\operatorname{fib}(\id_{\Kstab}\xrightarrow{\eta}iL),
 \qquad
 \bigcirc_X:=\Residual(X)=\operatorname{fib}(\eta_X:X\to iLX)
\]
と定める。安定性により自然なファイバー列
\[
 \bigcirc_X\longrightarrow X\xrightarrow{\eta_X}iLX
 \longrightarrow\Sigma\bigcirc_X
\]
が存在する。
\end{definition}

\begin{proposition}[外点の観測消去]
任意の$X\in\Kstab$に対し、$L(\bigcirc_X)\simeq0$である。
\end{proposition}
\begin{proof}
$L$は完全局所化であり、$L(\eta_X):LX\to L iLX$は同値である。
したがって$L$を上のファイバー列へ適用すると、そのファイバー
$L(\bigcirc_X)$は零対象と同値になる。
\end{proof}

\begin{definition}[$S_4$領域]
体系内不可記述領域を、観測局所化の核
\[
 \KSfour:=\ker L
 =\{X\in\Kstab\mid LX\simeq0\}
\]
として定義する。これは「体系の外」ではなく、五蘊圏内部で定義される充満部分圏である。
\end{definition}

\begin{axiom}[外点生成公理]
$\KSfour$は非零であり、写像空間関手
\[
 \operatorname{Map}_{\Kstab}(\bigcirc,-):\KSfour\longrightarrow\mathcal{S}
\]
が同値を検出する生成対象$\bigcirc\in\KSfour$を持つ。
単一生成対象が存在しないモデルでは、生成族$\{\bigcirc_\alpha\}$を許す。
\end{axiom}

\begin{remark}[$S_3$と$S_4$の関係]
$S_3$の観測可能成分$iLX$と$S_4$の不可視残余$\bigcirc_X$は、
$S_3\to S_4$という単純な逐次写像で結ばれるのではない。同一の$X$から
ファイバー列によって同時に分解される相補的成分である。五句循環における
「$S_3\to S_4$」は、観測結果$LX$だけでは$X$を復元できないことから
$\bigcirc_X$を検出する操作として読む。
\end{remark}

\subsection{強制法的外点の位置づけ}

可算推移モデル$M$、強制半順序$\mathbb{P}\in M$、$M$-ジェネリック
フィルター$G$を用いると、標準名$\dot G\in M$とその解釈
$\dot G^G=G\notin M$は、基底世界に存在しない対象の意味論的候補を与える
\cite{Jech2003}。
ただし、これを外点の具体例$\bigcirc_G$とするには、対応する安定圏$\Kstab(M,\mathbb P)$と
局所化$L$を構成し、$L(\bigcirc_G)\simeq0$を別途証明しなければならない。
したがって強制法は現段階では外点の\textbf{候補モデル}であり、
$\LawvereInfinity$の定義または存在証明には用いない。

%=============================================================
\section{五蘊モナドの圏論的定義}
%=============================================================

\subsection{色蘊：搾取モナド \texorpdfstring{$\Mdelta$}{M-delta}}

\begin{definition}[搾取モナド]
圏 $\mathcal{C}$（粒度差のない基底圏）と
圏 $\mathcal{D}$（$\Mdelta$ 文脈を持つ制度圏）の間に随伴
\begin{equation}
F : \mathcal{C} \rightleftharpoons \mathcal{D} : G, \quad F \dashv G
\end{equation}
を設定し、搾取モナドを
\begin{equation}
\Mdelta = G \circ F : \mathcal{C} \to \mathcal{C}
\end{equation}
と定義する。

単位 $\eta^\delta : \id_{\mathcal{C}} \Rightarrow \Mdelta$ は
価値・主体を制度文脈に埋め込む射である。
乗法 $\mu^\delta : \Mdelta\Mdelta \Rightarrow \Mdelta$ は
$G\varepsilon F$ により与えられ、$\varepsilon : FG \Rightarrow \id_{\mathcal{D}}$ は
随伴の余単位である。
\end{definition}

\begin{proposition}[搾取モナドのモナド則]
$(\Mdelta, \eta^\delta, \mu^\delta)$ はモナド則を満たす。
\end{proposition}
\begin{proof}
随伴 $F \dashv G$ から生成されるモナドは一般にモナド則を満たす
（Mac Lane \cite{MacLane1971}、VI章定理1）。
\end{proof}

\subsection{受蘊：二重否定モナド \texorpdfstring{$\Mneg$}{M-negneg}}

\begin{definition}[二重否定モナド]
トポス $\mathcal{E}$（あるいは五蘊圏の基礎圏 $\mathcal{C}$）において、
局所作用素 $j = \neg\neg : \Omega \to \Omega$ に対応する
層化随伴
\begin{equation}
a : \mathcal{E} \rightleftharpoons \mathrm{Sh}_{\neg\neg}(\mathcal{E}) : i,
\quad a \dashv i
\end{equation}
から生成されるモナドを二重否定モナドとする。
\begin{equation}
\Mneg = i \circ a : \mathcal{E} \to \mathcal{E}
\end{equation}
\end{definition}

\begin{theorem}[冪等性]
二重否定モナドは冪等である。すなわち、乗法
$\mu^\neg : \Mneg\Mneg \Rightarrow \Mneg$ は同型射である。
\end{theorem}
\begin{proof}
層化関手 $a$ は冪等である。$\mathrm{Sh}_{\neg\neg}(\mathcal{E})$ の対象は
すでに $\neg\neg$-層であるため、$a(iF) \cong F$ が成立する。
したがって $\Mneg \circ \Mneg \cong \Mneg$ となる。
\end{proof}

\subsection{想蘊：洗練モナド \texorpdfstring{$\Mr$}{M-r}}

\begin{definition}[洗練モナド]
圏 $\mathcal{C}$ 上の自己関手 $\Mr : \mathcal{C} \to \mathcal{C}$ と
自然変換 $\eta^r : \id \Rightarrow \Mr$、$\mu^r : \Mr\Mr \Rightarrow \Mr$
からなる三つ組 $(\Mr, \eta^r, \mu^r)$ を洗練モナドとする。

$\Mr$ は対象 $A$ を「洗練過程の中にある $A$」へ持ち上げる。
単位 $\eta^r_A : A \to \Mr A$ は文脈化を表し、
乗法 $\mu^r_A : \Mr\Mr A \to \Mr A$ は洗練の洗練を一段階へ統合する。
\end{definition}

\begin{proposition}[洗練順序との整合性]
$\Mr$ は洗練順序 $\sqsupseteq$ に関して単調である。
すなわち、$a \sqsupseteq a'$ ならば $\Mr a \sqsupseteq \Mr a'$。
\end{proposition}

\subsection{行蘊：カット除去モナド \texorpdfstring{$\Mc$}{C}}

\begin{definition}[カット除去モナド]
証明図の圏 $\mathcal{P}$ 上の自己関手 $\Mc : \mathcal{P} \to \mathcal{P}$ を、
対象 $P$ に対してカット規則を除去した証明図 $\Mc(P) = \mathrm{CutElim}(P)$
へ写す関手として定義する。

単位 $\eta^c : \id \Rightarrow \Mc$ は「カット規則を適用しない」埋め込みであり、
乗法 $\mu^c : \Mc\Mc \Rightarrow \Mc$ は多重カットの統合である。
\end{definition}

\begin{theorem}[カット除去モナドのモナド則]
$(\Mc, \eta^c, \mu^c)$ はモナド則を満たす。
\end{theorem}
\begin{proof}
カット消去アルゴリズムの正規化可能性に基づく。
カット消去済みの証明図にさらにカット消去を適用しても変化しないことから
単位元則が、カット消去の反復が結合的であることから結合則が従う。
\end{proof}

\subsection{識蘊：統制モナド \texorpdfstring{$\T$}{T}}

\begin{definition}[統制モナド]
統制モナド $\T$ は、五蘊圏 $\Sk$ において
以下の固定点方程式で定義される。
\begin{equation}
\T \cong F(\T), \qquad F(X) = X(\Mdelta, \Mneg, \Mr, \Mc, X)
\end{equation}
ここで $F$ は4つのモナドと $X$ 自身を引数に取る関手である。
\end{definition}

\begin{remark}
$\T$ が自身を包含することは、識（第5蘊）が全ての蘊を統括しながら
自身も空であるという仏教的構造と対応する。
\end{remark}

%=============================================================
\section{統制構造と識軸スパイラル}
%=============================================================

\subsection{再帰的構成によるフロベニウス性}

\begin{theorem}[再帰的構成によるフロベニウス性]
固定点 $\T \cong F(\T)$ において、以下の二方向が成立する。
\begin{align}
\fold^\T &: F(\T) \to \T \quad \text{（前進：折り畳み）} \\
\unfold^\T &: \T \to F(\T) \quad \text{（後退：展開）}
\end{align}
固定点方程式がこの二方向を等式として同一視するため、
フロベニウス条件
\begin{equation}
(\fold^\T \circ \T\unfold^\T) \cong (\unfold^\T \circ \fold^\T)
\end{equation}
が固定点方程式の帰結として成立する。
\end{theorem}
\begin{proof}
固定点 $\T \cong F(\T)$ は同型射 $\alpha : F(\T) \xrightarrow{\sim} \T$ および
その逆 $\alpha^{-1} : \T \xrightarrow{\sim} F(\T)$ を与える。
$\fold^\T := \alpha$（折り畳み）、$\unfold^\T := \alpha^{-1}$（展開）と定めると、
$\fold^\T$ と $\unfold^\T$ は互いに逆射であるため
\begin{equation}
\fold^\T \circ \unfold^\T = \id_{\T}, \qquad
\unfold^\T \circ \fold^\T = \id_{F(\T)}
\end{equation}
が成立する。これよりフロベニウス条件が $\id$ の可換性として得られる。
\end{proof}

\begin{remark}
通常のモナドは前進方向のみを持つ。
しかし統制モナドの固定点方程式は自然に二方向を生成する。
これが五蘊圏におけるフロベニウス性の根源である。
\end{remark}

\subsection{識軸スパイラル}

\begin{definition}[識軸スパイラル]
統制モナド $\T$ による軸的作用を
\begin{equation}
\SpiralT : (\Mdelta, \Mneg, \Mr, \Mc) \longrightarrow 
             (\Mneg, \Mr, \Mc, \Mdelta)
\end{equation}
と定める。この反復 $\SpiralT^n$ が四句分別の巡回を与え、
$\SpiralT^4(A)$の観測可能成分が$A$の観測可能成分へ戻るとき、五蘊圏は
\textbf{観測的に閉じた}識軸スパイラルを持つという。
\[
 L\SpiralT^4(A)\simeq LA.
\]
この閉路の不可視なずれは
\[
 \bigcirc_A:=\Residual(\SpiralT^4(A))
\]
によって測られる。
\end{definition}

\begin{proposition}[スパイラルの残余閉包条件]
観測的閉包$L\SpiralT^4(A)\simeq LA$のもとで、$\bigcirc_A$は
一周後に観測から失われる情報を完全に測る。$\bigcirc_A\simeq0$ならば閉路は
安定実現においても閉じる。$\bigcirc_A\not\simeq0$ならば、閉路は$S_3$では閉じるが
$S_4$に非自明な残余を残す。
\end{proposition}

\begin{remark}
スパイラルだけから後退方向やフロベニウス構造が自動的に生成されるわけではない。
そのためには残余列の分裂、随伴、または前進・後退を比較する追加自然変換が必要である。
$\bigcirc_A$は、その追加構造が失敗する障害を記録する。
\end{remark}

\begin{remark}
この $\SpiralT$ は、Hopf代数における反極（antipode）に類似する役割を持つ。
すなわち、代数的な前進方向と余代数的な後退方向を同一対象上で接続し、
方向反転を媒介する操作である。
\end{remark}

%=============================================================
\section{証明論的固定点\texorpdfstring{$\LawvereInfinity$}{Lambda-infinity}との比較}
%=============================================================

本節は外点$\bigcirc$を定義するものではない。$\LawvereInfinity$は
証明論的自己言及の独立した固定点であり、その詳細は別稿
\cite{Chiba2026limit}に委ねる。以下では統制固定点との構造的類似のみを記録する。

\subsection{自己言及演算子としての識蘊}

\begin{definition}[自己言及演算子 $\SigmaSk$]\label{def:self-reference-op}
五蘊圏 $\Sk$ の言語における文 $A$ に対し、自己言及演算子 $\SigmaSk$ を以下で定義する。
\begin{equation}
\SigmaSk(A) := \lceil A \text{ は } \Sk \text{ において証明不可能である} \rceil 
= \neg\PrSk(\ulcorner A \urcorner)
\end{equation}
$\SigmaSk(A)$の算術的構成には対角補題を用いる\cite{Chiba2026limit}。
\end{definition}

\begin{definition}[自己言及の反復列]\label{def:self-reference-seq}
文 $A_0$ に対し、$\SigmaSk$ の反復列を以下で定義する。
\begin{equation}
A_0, \quad A_1 = \SigmaSk(A_0), \quad A_2 = \SigmaSk(A_1), \quad \ldots, \quad A_n = \SigmaSk^n(A_0)
\end{equation}
\end{definition}

\begin{definition}[Lawvere無限階層の形式体系的実現]\label{def:e-infinity}
自己言及演算子$\SigmaSk$を$S_4$の作用素$\Sigma_4$として実現する。
以下の等式を満たす文は、Lawvere無限階層$\LawvereInfinity$の各周回における
形式体系的実現である。
\begin{equation}
\LawvereInfinity = \SigmaSk(\LawvereInfinity)
\end{equation}
すなわち $\LawvereInfinity$ は「$\LawvereInfinity$ は $\Sk$ において証明不可能である」という内容の文であり、
\begin{equation}
\Sk \vdash \LawvereInfinity \leftrightarrow \neg\PrSk(\ulcorner \LawvereInfinity \urcorner)
\end{equation}
を満たす。
この証明論的固定点は、外点$\bigcirc$または$\KSfour$の生成対象とは同一視しない。
\end{definition}

\begin{remark}[$=$ の正当性：不動点の静的記述として]\label{rem:equality-justification}
定義 \ref{def:e-infinity} における $=$ は、対象の自己同一性を要請する $A=A$ の $=$ とは種類が異なる。
$\lim_{n\to\infty} a_n = L$ において $=$ が無限プロセスを静的値に変換するように、
$\LawvereInfinity = \SigmaSk(\LawvereInfinity)$ はプロセス $\SigmaSk$ の反復が必然的に到達する不動点を静的値として記述する。
これはバナッハの不動点定理における $x = f(x)$ と同型の操作であり、プロセスの収束先に名前をつける発見的記述である。
「極限」という概念が本質的に静的なトリックである以上、この $=$ は正当に使用できる。
\end{remark}

\subsection{真理保存的実現における固定点}

\begin{theorem}[真理保存的実現における$\LawvereInfinity$]\label{thm:e-infinity-existence}
五蘊圏の基礎圏を十分強い真理保存体系の圏として実現した場合、
\begin{equation}
\exists \LawvereInfinity \quad \text{s.t.} \quad 
\Sk \vdash \LawvereInfinity \leftrightarrow \neg\PrSk(\ulcorner \LawvereInfinity \urcorner)
\end{equation}
が成立する。これは証明論的実現についての主張であり、観測局所化の核に属する
$\bigcirc$の存在とは独立である。
\end{theorem}
\begin{proof}
対角補題を論理式$\varphi(x):=\neg\PrSk(x)$に適用する。すると、
\begin{equation}
\Sk \vdash G \leftrightarrow \neg\PrSk(\ulcorner G \urcorner)
\end{equation}
を満たす文 $G$ が $\Sk$ の内部で構成可能である。この $G$ を $\LawvereInfinity$ と命名する。

$\LawvereInfinity$ の構成は $\Sk$ の公理・推論規則・ゲーデル符号化のみを使用し、外部からの付加を一切必要としない。したがって $\LawvereInfinity$ は $\Sk$ の言語で表現可能な文として $\Sk$ に内在する。
\end{proof}

\begin{remark}[「発見」としての $\LawvereInfinity$]\label{rem:e-infinity-discovery}
ゲーデルはこの $G = \LawvereInfinity$ を「$\Sk$ において証明も反証もできない文」として提示し、体系の不完全性の証拠とした。
しかし定理 \ref{thm:e-infinity-existence} が示すのは、$\LawvereInfinity$ が $\Sk$ の外部にあるのではなく、$\Sk$ 自身の対角補題から必然的に生成される内在的対象であるということである。
「決定不能」として現れたのは、$\Sk$ の真理値空間 $\{0,1\}$（$\mathcal{B}$：二値截断部分圏）が $\LawvereInfinity$ に対応する値を持たなかっただけである。

リーマン球面において $\infty$ が「複素平面には存在しないが、リーマン球面上では通常の点として存在する」のと同型の構造である。
数学社会は $\mathcal{B}$ を唯一の真理値空間として採用し続けたために、$\LawvereInfinity$ を「壁」として誤読してきた。これは発明ではなく、未発見だったものの発見である。
\end{remark}

\begin{proposition}[統制モナドと対角補題の対応]\label{prop:monad-diagonal}
基礎圏 $\mathcal{C}$ を真理保存体系の圏 $\mathbf{Th}$ とする。
各対象 $T \in \mathbf{Th}$（形式体系）に対し、
自己言及演算子 $\Sigma_T$ は $\mathcal{C}$ 上の自己関手として作用し、
その不動点 $\LawvereInfinity$ は統制モナド $\T_T$ の固定点
$\T_T \cong F(\T_T)$ の具体例として現れる。

ただし、この対応は統制モナドと自己言及固定点の間のものであり、
$\LawvereInfinity$と外点$\bigcirc$の対応を意味しない。
\end{proposition}

\begin{remark}
$\LawvereInfinity$ は統制モナドの固定点の「部分例」に過ぎない。
統制モナドは任意の圏 $\mathcal{C}$ 上の関手の固定点として定義され、より広い普遍性を持つ。
正しい位置づけは：
\begin{quote}
\textbf{五蘊圏の統制モナドは「自己言及固定点」の普遍構造であり、$\LawvereInfinity$ はその構造が形式体系論において実現された具体例である。}
\end{quote}
$\LawvereInfinity$ は五蘊圏を\textbf{還元}するのではなく、\textbf{補強}する。
\end{remark}

\subsection{\texorpdfstring{$\LawvereInfinity$}{Lambda-infinity}の否定と二諦の非閉鎖性}

\begin{definition}[$\neg\LawvereInfinity$ の定義]\label{def:neg-e-infinity}
$\LawvereInfinity$ の否定を以下で定義する。
\begin{equation}
\neg\LawvereInfinity = \PrSk(\ulcorner \LawvereInfinity \urcorner)
\end{equation}
すなわち「$\LawvereInfinity$ は $\Sk$ において証明可能である」。
\end{definition}

\begin{theorem}[$\LawvereInfinity$ と $\neg\LawvereInfinity$ の非閉鎖性]\label{thm:non-closure}
$\neg\LawvereInfinity \neq \LawvereInfinity$。
\end{theorem}
\begin{proof}
定義 \ref{def:neg-e-infinity} より $\LawvereInfinity = \neg\PrSk(\ulcorner \LawvereInfinity \urcorner)$ であるのに対し、$\neg\LawvereInfinity = \PrSk(\ulcorner \LawvereInfinity \urcorner)$ である。両者は構文的に異なる文である。
\end{proof}

\begin{remark}[フロベニウス条件からの独立性]
$\LawvereInfinity$の非閉鎖性は証明論的な構文上の性質であり、五蘊圏の
フロベニウス条件を導出も補強もしない。両者を接続するには、証明論的圏から
五蘊圏の安定実現への関手と、その関手が固定点および前進・後退比較を保存することを
別途示す必要がある。本稿ではそのような関手を仮定しない。
\end{remark}

\begin{remark}
これは、Priest の「パラドックスの論理」\cite{Priest1979} に代表される矛盾許容論理との構造的差異である。
同論理では、真かつ偽である値が否定に関して固定的に扱われ、矛盾が推論空間内の安定した領域として保持される。

これに対して、$\LawvereInfinity$ と $\neg\LawvereInfinity$ は互いに異なる文として、$\SigmaSk$ のサイクル上の二点をなす。したがって本体系において、矛盾は閉じた終端ではなく通過点として機能する。
\end{remark}

\subsection{真理保存五蘊圏の不完全性定理}

\begin{theorem}[五蘊圏の不完全性定理]
真理保存五蘊圏 $\Sk$ において、$\LawvereInfinity$ は体系内部から証明も反証もできない。
すなわち、
\begin{equation}
\Sk \nvdash \LawvereInfinity \quad \text{かつ} \quad \Sk \nvdash \neg\LawvereInfinity.
\end{equation}
\end{theorem}
\begin{proof}
$\Sk$の証明論的実現が一貫していると仮定する。
\begin{enumerate}[label=(\roman*)]
\item $\Sk \vdash \LawvereInfinity$ と仮定すると、$\Sk \vdash \neg\PrSk(\ulcorner \LawvereInfinity \urcorner)$ となるが、
  これは「$\LawvereInfinity$ は証明可能である」という事実と矛盾する。
\item $\Sk \vdash \neg\LawvereInfinity$ と仮定すると、$\Sk \vdash \PrSk(\ulcorner \LawvereInfinity \urcorner)$ となるが、
  これは $\Sk$ の一貫性より $\Sk \vdash \LawvereInfinity$ を導き、再び矛盾する。
\end{enumerate}
したがって両方とも証明不可能である。
\end{proof}

\begin{remark}
この定理は証明論的固定点の存在を述べるものであり、外点$\bigcirc$の存在根拠ではない。
$\bigcirc$は観測局所化の核によって、$\LawvereInfinity$は対角補題によって、
それぞれ独立に定義される。
\end{remark}

%=============================================================
\section{五蘊圏の高次構造}
%=============================================================

\subsection{2圏論的解釈}

五蘊圏は、モナドの2圏 $\mathbf{Mnd}(\mathcal{C})$ 上の構造として理解できる。

\begin{definition}[五蘊2圏]
圏 $\mathcal{C}$ 上のモナドの2圏 $\mathbf{Mnd}(\mathcal{C})$ は、
以下を持つ2圏である。

\begin{itemize}
\item \textbf{0-胞体（対象）}：圏 $\mathcal{C}$
\item \textbf{1-胞体（1-射）}：$\mathcal{C}$ 上のモナド
  $\Mdelta, \Mneg, \Mr, \Mc, \T$
\item \textbf{2-胞体（2-射）}：モナド間の自然変換（分配律・モナド射）
\end{itemize}
\end{definition}

\begin{proposition}[五蘊2圏における統制]
統制モナド $\T$ は、五蘊2圏 $\mathbf{Mnd}(\mathcal{C})$ において
\textbf{終対象（terminal object）}として振る舞う。
すなわち、任意のモナド $\mathcal{M} \in \{\Mdelta, \Mneg, \Mr, \Mc\}$ に対し、
唯一のモナド射 $\mathcal{M} \to \T$ が存在する。
\end{proposition}

\begin{remark}
五蘊2圏において、4つの従属モナドは1-射として配置され、
統制モナドはそれらの「極限」として位置づけられる。
識軸スパイラルは、この2圏内の\textbf{巡回自己2-射}として定式化される。
\end{remark}

\subsection{ホップf代数構造}

\begin{definition}[五蘊型ホップf代数]
五蘊圏 $\Sk$ は、以下の構造を持つ\textbf{五蘊型ホップf代数}として理解される。

\begin{itemize}
\item \textbf{代数構造}（乗法・単位）：
  前進方向のモナド演算 $\fold^\T : F(\T) \to \T$
\item \textbf{余代数構造}（余乗法・余単位）：
  後退方向の余モナド演算 $\unfold^\T : \T \to F(\T)$
\item \textbf{反極 $S$}：識軸スパイラル $\SpiralT$。
  これは方向反転を媒介する。
\end{itemize}
\end{definition}

\begin{proposition}[フロベニウス性とホップf代数]
五蘊型ホップf代数において、フロベニウス条件 \eqref{eq:frobenius} は
ホップf代数の公理
\begin{equation}
\fold^\T \circ (S \otimes \id) \circ \Delta 
= \eta \circ \varepsilon 
= \fold^\T \circ (\id \otimes S) \circ \Delta
\end{equation}
の圏論的類似として解釈される。
\end{proposition}

\begin{remark}
本稿における $\SpiralT$ は、この等式を厳密に仮定するものではなく、
識軸が前進方向と後退方向を接続することを表す構造的類比である。
\end{remark}

\subsection{トレース付きモノイダル圏}

\begin{proposition}[トレース構造]
識軸スパイラルの反復は、トレース付きモノイダル圏における
\textbf{フィードバック操作}に対応する。

射 $f : A \otimes U \to B \otimes U$ に対するトレース
\begin{equation}
\trace^U_{A,B}(f) : A \to B
\end{equation}
において、$U$ を識軸とみなすと、一回のトレースが
四句分別の一巡を閉じる操作に相当する。
したがって、識軸スパイラルは、反極的な方向反転と
トレース的な再帰の双方を持つ。
\end{proposition}

%=============================================================
\section{二諦の公理と外在・内在}
%=============================================================

\subsection{二重五蘊構造}

\begin{definition}[外在・内在の五蘊]
五蘊圏 $\Sk$ は、外在と内在の二重構造を持つ。
\begin{align}
\TExt &= (\Mdelta^{\ext}, \Mneg^{\ext}, \Mr^{\ext}, \Mc^{\ext}, \T^{\ext}) \\
\TInt &= (\Mdelta^{\intn}, \Mneg^{\intn}, \Mr^{\intn}, \Mc^{\intn}, \T^{\intn})
\end{align}
\end{definition}

\begin{axiom}[相互完全観測不可能性]
外在の五蘊 $\TExt$ と内在の五蘊 $\TInt$ は相互に完全に観測不可能である。
\begin{equation}
\TExt \not\leftrightarrow \TInt, \qquad
\TInt \not\leftrightarrow \TExt \quad \text{（直接観測の不可能性）}
\end{equation}
ただし「完全に観測不可能」は「常に観測可能」と同義であり、
内在側は直接対象として観測できないが、
全ての外在的観測の地盤として常に作動している。
\end{axiom}

\begin{definition}[残余を介した二諦的連携]
外在と内在の連携は非対称的である。各対象$X$について、観測可能成分への単位
$\eta_X:X\to iLX$と、不可視残余のファイバー包含
\[
 m_X:\bigcirc_X\longrightarrow X
\]
を二方向の比較データとする。$\eta_X$は世俗諦への観測を、$m_X$は観測から
消去される勝義諦的残余が現象へ関与する仕方を表す。
\end{definition}

\begin{proposition}[二諦的連携の障害]
観測可能成分上の前進・後退比較が同値であっても、局所化前のフロベニウス比較が
同値になるとは限らない。その障害は$\bigcirc_X\in\KSfour$に属する。
$\bigcirc_X\simeq0$であるか、残余列の両立的な分裂が追加される場合に限り、
観測上の比較を局所化前へ持ち上げられる。
\end{proposition}
\begin{proof}
$L(\bigcirc_X)\simeq0$なので、$L$を適用した比較は残余を検出しない。
したがって観測上の同値から局所化前の同値を結論するには、核$\ker L$に残る
成分が零であること、またはその成分を整合的に分裂させる追加データが必要である。
\end{proof}

%=============================================================
\section{五蘊・四句分別・二諦との対応}
%=============================================================

\subsection{五蘊とモナドの対応}

\begin{proposition}[五蘊とモナドの対応]
五蘊圏の構成は、仏教の五蘊と以下の対応を持つ。

\begin{center}
\begin{tabular}{llll}
\hline
\textbf{五蘊} & \textbf{モナド} & \textbf{四句} & \textbf{方向性} \\
\hline
色（形・物質） & $\Mdelta$（搾取） & 是 & 前進のみ \\
受（感受） & $\Mneg$（二重否定） & 非 & 前進のみ \\
想（知覚・表象） & $\Mr$（洗練） & 亦是亦非 & 双方向（誘導） \\
行（意志・形成力） & $\Mc$（カット除去） & 非是非非 & 双方向（誘導） \\
識（識別・意識） & $\T$（統制） & 統制 & 再帰的 \\
\hline
\end{tabular}
\end{center}
\end{proposition}

\subsection{四句分別の圏論的解釈}

\begin{proposition}[四句分別の圏論的解釈]
四句分別は、五蘊圏におけるモナド間の関係の位相として現れる。

\begin{align}
\phi_S \text{（是）} &\longleftrightarrow \Mdelta 
\quad \text{（実現子の存在：} e \Vdash A \text{）} \\
\phi_N \text{（非）} &\longleftrightarrow \Mneg 
\quad \text{（実現子の不在：} \neg(e \Vdash A) \text{）} \\
\phi_{SN} \text{（亦是亦非）} &\longleftrightarrow \Mr 
\quad \text{（実現子列の収束：洗練過程）} \\
\phi_{NS} \text{（非是非非）} &\longleftrightarrow \Mc 
\quad \text{（残余の確定：カット除去後の残余）}
\end{align}
\end{proposition}

\subsection{二諦との対応}

\begin{proposition}[二諦との対応]
\begin{align}
\TExt \text{（外在の五蘊）} &\longleftrightarrow 
\text{世俗諦（観測可能な現象の層）} \\
\TInt \text{（内在の五蘊）} &\longleftrightarrow 
\text{勝義諦（空・中道不動点の層）}
\end{align}

世俗諦は$LX$として観測可能化され、勝義諦的残余は
$m_X:\bigcirc_X\to X$を通じて現象$X$へ関与する。この記述は
$\bigcirc_X$を「体系の外」に置かず、観測局所化の核に内在させる。
\end{proposition}

%=============================================================
\section{五蘊圏と既存理論の位置づけ}
%=============================================================

\subsection{包含関係}

五蘊圏 $\Sk$ は、以下の既存理論を包含する。

\begin{center}
\begin{tabular}{ll}
\hline
\textbf{既存理論} & \textbf{五蘊圏内の位置づけ} \\
\hline
モナドの2圏 $\mathbf{Mnd}(\mathcal{C})$ & 五蘊2圏の部分構造 \\
ホップf代数 & 五蘊型ホップf代数の代数・余代数部分 \\
トレース付きモノイダル圏 & 識軸スパイラルのトレース解釈 \\
実現子論的圏 $\Asm(\mathcal{K})$ & 基礎圏 $\mathcal{C}$ の具体例 \\
トポス論（層化随伴） & 二重否定モナドの基礎構造 \\
\hline
\end{tabular}
\end{center}

\subsection{五蘊圏の普遍性}

五蘊圏は、個別のモナドの内容（搾取の社会的構造、カット除去の証明論など）を
抽象化し、「5つのモナドが統制モナドを軸として相互作用する」という
普遍的パターンを記述する。

この普遍性は、五蘊圏が\textbf{圏の圏（category of categories）}における
\textbf{普遍構成（universal construction）}として位置づけられることを示唆する。

\subsection{\texorpdfstring{$\LawvereInfinity$}{Lambda-infinity}との関係}

\begin{proposition}[外点と$\LawvereInfinity$の型の分離]
外点$\bigcirc$は$\KSfour=\ker L$の対象または生成元であり、
$\LawvereInfinity$は形式体系の言語に属する自己言及文である。
したがって、両者の間に標準的な等式・同型・対応関手は存在せず、本稿では
$\bigcirc\cong\LawvereInfinity$を主張しない。

\begin{center}
\begin{tabularx}{\textwidth}{@{}l >{\raggedright\arraybackslash}X >{\raggedright\arraybackslash}X@{}}
\toprule
側面 & 外点$\bigcirc$ & $\LawvereInfinity$ \\
\midrule
型 & $\Kstab$の対象 & 形式体系の文 \\
普遍問題 & 観測局所化の不可視残余 & 自己言及作用素の固定点 \\
存在根拠 & $\ker L$の非自明性と生成公理 & 対角補題 \\
強制法 & 候補モデルを与えうる & 定義には用いない \\
\bottomrule
\end{tabularx}
\end{center}
\end{proposition}

\begin{remark}[将来の翻訳問題]
両者の接続を検討するには、例えば$\KSfour$から体系$T$の超教義または
構文圏$\mathcal{E}_T$への関手
\[
 \Phi:\KSfour\longrightarrow\mathcal{E}_T
\]
を構成し、$\Phi(\bigcirc)$が$\LawvereInfinity$を表現する条件を証明する必要がある。
これは現段階の定理ではなく、独立した研究課題である。
\end{remark}

%=============================================================
\section{結論}
%=============================================================

本稿は五蘊圏$\Sk$を公理化し、外点$\bigcirc$へ観測局所化による
圏論的定義を与えた。

\begin{enumerate}
\item \textbf{基本定義}：五蘊圏を5つのモナド・統制固定点・識軸スパイラル・
  フロベニウス条件・二諦の公理からなる圏論的構造として定義した。

\item \textbf{五蘊モナド}：搾取・二重否定・洗練・カット除去・統制の
  各モナドを随伴・層化・自己関手として定式化した。

\item \textbf{統制構造}：統制モナド $\T$ の固定点方程式と識軸スパイラル $\SpiralT$ により、
  観測的閉包と局所化前に残る障害を区別した。

\item \textbf{外点の定義}：安定実現$\Kstab$、観測局所化$L$、不可視残余
  $\bigcirc_X=\operatorname{fib}(X\to iLX)$を導入し、$S_4$領域を
  $\KSfour=\ker L$として定義した。

\item \textbf{高次構造}：五蘊圏を2圏論・ホップf代数・トレース付きモノイダル圏の
  交差点として位置づけた。

\item \textbf{五蘊対応}：五蘊・四句分別・二諦と圏論的構造の間の
  同型対応を明示した。

\item \textbf{$\LawvereInfinity$との分離}：$\LawvereInfinity$を証明論的固定点、
  $\bigcirc$を局所化の核の対象として区別した。両者を結ぶ翻訳関手の構成は
  将来課題として残した。
\end{enumerate}

五蘊圏は、個別のモナドの内容ではなく、それらを統合する\textbf{構造の構造}として機能する。
識（統制モナド）は、五蘊圏の中で自己言及的固定点をなし、
前進と後退・外在と内在・代数と余代数を統合する軸として働く。
この自己言及性は、圏論の核が本質的に再帰的であることを、
五蘊構成を通じて反映している。

外点$\bigcirc$は「体系の外」に置かれた存在ではなく、観測局所化によって
消去される内在的残余である。この定義により、観測可能性と存在を区別しながら、
二諦・スパイラル・フロベニウス比較に現れる不可視成分を同一の圏論的言語で扱える。
$\LawvereInfinity$との類似は自己言及と不可視性に関する発見的比較にとどまり、
数学的同一視には別途翻訳定理を要する。

\begin{thebibliography}{99}

\bibitem{Abramsky1991}
Samson Abramsky.
\newblock Domain theory in logical form.
\newblock \emph{Annals of Pure and Applied Logic}, 51(1--2):1--77, 1991.

\bibitem{Bauer2000}
Andrej Bauer.
\newblock The realizability approach to computable analysis and topology.
\newblock Ph.D. thesis, Carnegie Mellon University, 2000.

\bibitem{Birkedal2000}
Lars Birkedal.
\newblock Developing theories of types and computability via realizability.
\newblock Ph.D. thesis, University of Edinburgh, 2000.

\bibitem{LurieHA}
Jacob Lurie.
\newblock \emph{Higher Algebra}.
\newblock 2017.

\bibitem{Jech2003}
Thomas Jech.
\newblock \emph{Set Theory}.
\newblock Springer, 3rd millennium edition, 2003.

\bibitem{Escardo2004}
Mart\'{\i}n Escard\'{o}.
\newblock Synthetic topology of data types and classical spaces.
\newblock \emph{Electronic Notes in Theoretical Computer Science}, 87:21--156, 2004.

\bibitem{Hoare1998}
C.~A.~R. Hoare and He Jifeng.
\newblock \emph{Unifying Theories of Programming}.
\newblock Prentice Hall, 1998.

\bibitem{Kleene1952}
Stephen Cole Kleene.
\newblock \emph{Introduction to Metamathematics}.
\newblock North-Holland, 1952.

\bibitem{MacLane1971}
Saunders Mac Lane.
\newblock \emph{Categories for the Working Mathematician}.
\newblock Springer, 1971.

\bibitem{Moggi1991}
Eugenio Moggi.
\newblock Notions of computation and monads.
\newblock \emph{Information and Computation}, 93(1):55--92, 1991.

\bibitem{Street2004}
Ross Street.
\newblock Frobenius monads and pseudomonoids.
\newblock \emph{Journal of Mathematical Physics}, 45(10):3930--3948, 2004.

\bibitem{JoyalStreetVerity1996}
Andr\'{e} Joyal, Ross Street, and Dominic Verity.
\newblock Traced monoidal categories.
\newblock \emph{Mathematical Proceedings of the Cambridge Philosophical Society},
  119(3):447--468, 1996.

\bibitem{Sweedler1969}
Moss E. Sweedler.
\newblock \emph{Hopf Algebras}.
\newblock W. A. Benjamin, 1969.

\bibitem{Priest1979}
Graham Priest.
\newblock The logic of paradox.
\newblock \emph{Journal of Philosophical Logic}, 8(1):219--241, 1979.

\bibitem{Chiba2026limit}
千葉将臣.
\newblock 真理保存性を持つ体系が必然的に $\Lambda_\infty$ を含む：対角補題によるLawvere無限階層の内在的定式化.
\newblock 2026.

\bibitem{Chiba2026cutmonad}
千葉将臣.
\newblock カット消去モナド：証明構造の安定化を担うモナド的定式化.
\newblock 未刊行原稿, 2026.

\bibitem{Chiba2026continuation}
千葉将臣.
\newblock 継続残余としての証明障害：実行可能性と静的証明のあいだの初等モデル.
\newblock 未刊行原稿, 2026.

\end{thebibliography}

\end{document}
